%=============================================================================
\section{BVS as a technological decision-support system}
\label{sec:system}
%=============================================================================

BVS is implemented as a system with three layers. Separating them matters for
governance, because each layer has a different epistemic role and a different
failure mode.

\paragraph{Evidence layer.}
Published studies and reports are ingested as PDF files. A large language model
extracts, for each source, the reported sample sizes, response distributions,
and regression coefficients, and proposes a mapping from each source item to a
target construct and response scale. Every proposed mapping is exported to a
review table (source item, source scale, extracted statistic, transformation,
resulting pseudo-counts) so that it can be verified manually. The language
model is therefore an extraction aid whose output is auditable line by line; it
is not a source of evidence. One of the authors reviewed all 80 extracted entries.

\paragraph{Inference engine.}
A Bayesian network encodes the dependencies among demographic, attitudinal, and
behavioural variables. Its structure is learned from national microdata by
Markov chain Monte Carlo over graphs and its parameters from national counts;
an outcome absent from the national data is attached through a transfer model
with explicit assumptions, which later Spanish data update by conjugate fusion
(Section~\ref{sec:framework}). All acceptance
probabilities and intervals reported in this paper are computed by this layer.

\paragraph{Access layer.}
A conversational interface (ASAP~v6) makes the engine usable without
statistical training. A free-text query is parsed by a language model into a
structured profile (single profile, comparison of profiles, or overview); the
profile is passed to the inference engine, which returns the projected
probability and its 95\% interval; the result is
displayed as a probability card with a short narrative
(Figure~\ref{fig:gradio}). The language model translates in both directions and
generates no probabilities. The prototype offers 13 product categories. In this
paper only food-packaging acceptance is supported by the transfer model of
Section~\ref{sec:step2}. The other twelve categories use base rates taken from
the shares of the Finnish survey that judged recycled plastic suitable for each
product type \citep{ruokamo2022}, multiplied by 0.85, together with the
food-packaging coefficients scaled by a fixed factor between 0.4 and 0.7. They
are exploratory, are not estimated from data, and are not reported here.
Response latency and the accuracy of query parsing have not been evaluated and
are left to future work.

\begin{figure}[htbp]
\centering
\includegraphics[width=0.9\textwidth]{fig_gradio.png}
\caption{Query flow of the ASAP~v6 interface, with the output of the current
model for the most receptive profile of Section~\ref{sec:spread}.}
\label{fig:gradio}
\end{figure}

Table~\ref{tab:requirements} relates the layers to the requirements of
Section~\ref{sec:requirements}. The technological contribution lies in their
combination: evidence integration, scenario exploration, and later updating
become available within one traceable workflow, to users who could not run the
underlying model themselves.

\begin{table}[htbp]
\centering
\caption{Design requirements and the BVS components that address them.}
\label{tab:requirements}
\small
\begin{tabularx}{\textwidth}{lXX}
\toprule
\textbf{Requirement} & \textbf{Component} & \textbf{Limit} \\
\midrule
R1 Pre-observational & External-evidence priors and transfer mapping for
  outcomes absent from national data & Levels and effects are transferred from
  other populations \\
R2 Segment-level & Profile projections with Monte Carlo intervals; segment
  projections at the central specification & Intervals omit transfer and
  extraction uncertainty \\
R3 Updatable & Conjugate update of a Beta prior matched to the projection &
  Implemented, not yet applied to Spanish data \\
R4 Accessible & Natural-language interface; review tables and versioned
  outputs & Query parsing and user interpretation not yet evaluated \\
\bottomrule
\end{tabularx}
\end{table}

%=============================================================================
\section{Data and evidence provenance}
\label{sec:data}
%=============================================================================

\paragraph{National microdata.}
The primary dataset is CIS Study~3391 (ISSP Environment~IV, Spain, 2023;
$n = 2{,}254$), comprising 194 raw variables filtered to 90 substantive items on
demographics, environmental concern, institutional trust, willingness to pay,
and recycling behaviour \citep{cis2023}. It is used to learn the dependency
structure and to estimate the conditional distributions of observed variables.
The final weight PESOFIN (range 0.53--3.59; Kish design effect 1.05,
effective $n = 2{,}153$) is applied to the counts of the conditional
probability tables and to all population estimates. Structure is learned
without weights, because the BDeu score is not a well-defined marginal
likelihood for fractional counts.

\paragraph{Missing values.}
Missing values affect 31 of the 90 network variables: 3.9\% of all cells of the
90 variables, 11.2\% on average among the 31 affected, and at most 62\%. They
are treated differently in three places. (i)~For structure learning and the
conditional tables, they are imputed with the modal category. Available-case
analysis is not used here because, with BDeu, each parent set would be scored
on a different subsample and the scores of different graphs would cease to be
comparable within MC$^3$. Mode imputation inflates the modal category of the
variables with most missing values (Appendix~\ref{app:missing}). (ii)~The
transfer model uses complete cases on its five predictors ($n = 2{,}133$ of
$2{,}254$). (iii)~Validation metrics compare synthetic and CIS distributions
over the non-missing responses of each variable.

\paragraph{Target outcomes.}
CIS~3391 contains no measure of five outcomes relevant to implementation:
general acceptance of recycled plastics, acceptance in food packaging,
perceived safety, perceived quality, and willingness to pay a premium. The
statement is limited to this dataset; we have not established that no Spanish
survey measures them.

\paragraph{External evidence.}
The extraction corpus comprises 11 journal articles reporting aggregate
statistics (Appendix~\ref{app:studies}, Table~\ref{tab:studies}): about
3{,}200 independent respondents, counting each study once at its largest
sample, and 80 extracted item-level entries. Three levels must be kept apart.
\emph{Corpus}: all 11 studies and 80 entries. \emph{Retained evidence}: entries
that measure one of the five target outcomes and pass the review rules
(\texttt{evidence\_rules.csv}); they come from four studies
(Table~\ref{tab:i2_summary}). \emph{Transfer model}: the base rate and
coefficients for food-packaging acceptance, all from \citet{ruokamo2022}.
Entries for variables the CIS already measures were extracted but are not
used, and no external study influences the learned structure or the
generation of observed variables.

\paragraph{External benchmark.}
The Ipsos Global Advisor survey on plastic packaging (\citealp{ipsos2022};
Spain, $n = 1{,}000$) is reserved for external comparison. It contributes no
entry to the priors. A scoping review \citep{polyportis2022} was used to define
the construct taxonomy and contributes no data.

